\section{Capítulo~\ref{sec:chemoutput}. A saída química}

Os dois fios para o coração e suas velocidades diferentes, a realimentação negativa nas cascatas hormonais, o código pela frequência das porções, o gerador diário e seu ajuste pela luz foram medidos diretamente; o limite $K<8$ é um teorema da teoria de controle, deduzido no capítulo; as pulsações geradas pela própria realimentação da cascata são uma hipótese com um experimento forte a seu favor.

{\footnotesize
\setlength{\tabcolsep}{3pt}\renewcommand{\arraystretch}{1.15}
\begin{longtable}{>{\raggedright}p{0.5cm}>{\raggedright}p{4.0cm}>{\raggedright}p{5.6cm}>{\raggedright}p{2.3cm}>{\raggedright\arraybackslash}p{1.7cm}}
\toprule
N.º & Proposição & Experimento e o que foi medido & Fonte & Status \\
\midrule\endfirsthead
\toprule
N.º & Proposição & Experimento e o que foi medido & Fonte & Status \\
\midrule\endhead
1 & O marca-passo do coração bate sozinho cerca de $100$ vezes por minuto; em repouso o freio está pisado, e a frequência costuma ficar em torno de $60$--$70$ & Seres humanos, desligamento completo dos dois fios para o coração: a frequência própria é maior que a de repouso e diminui com a idade; em adultos, da ordem de $100$ batimentos por minuto. Logo, em repouso predomina o freio. A frequência de repouso de $60$--$70$ é um valor típico, sem citação à parte. & \cite{JoseCollison1970ihr} & medido \\
2 & O acelerador \nt{NORA} e o freio \nt{ACET} são filtros passa-baixas, e o freio é mais rápido~\eqref{eq:gasbrake} & Cão, estimulação modulada em frequência do nervo vago ou do simpático cardíaco e função de transferência até a frequência atrial: as duas vias são filtros passa-baixas; a simpática tem frequência de corte muito menor e um atraso puro de cerca de $1.7\,\text{s}$. As características dependem do tônus médio, de modo que a fórmula linear~\eqref{eq:gasbrake} é uma aproximação. & \cite{Berger1989} & medido \\
3 & O freio é rápido porque o \nt{ACET} abre o canal de potássio sem segundo mensageiro, e o \nt{NORA} age por uma cadeia de reações & Fragmentos de membrana de células atriais: o canal de potássio aberto pela acetilcolina é aberto pelas subunidades $\beta\gamma$ da proteína G ligada ao receptor, sem segundo mensageiro dentro da célula. A via do \nt{NORA} pelo AMPc não é citada aqui. & \cite{Logothetis1987gbg} & medido \\
4 & O sangue dá a volta no corpo em cerca de um minuto; o \nt{ADRE} da corrente sanguínea chega a todos os órgãos com receptor & Estimativa geral de ordem de grandeza; assim está formulada no capítulo, sem fonte citada. & --- & estimativa \\
5 & A cascata hormonal é envolvida por realimentação negativa: o hormônio final inibe os primeiros estágios & Revisão de experimentos: os hormônios do córtex adrenal suprimem a liberação de ACTH pela hipófise e do hormônio liberador pelo hipotálamo em três faixas de tempo: rápida, intermediária e lenta. & \cite{KellerWood1984fb} & medido \\
6 & Três estágios iguais são estáveis para $K<8$~\eqref{eq:cascadestable}; no limite, o período é $2\pi\tau/\sqrt3\approx3.6\tau$ & Deduzido no capítulo: na frequência $\sqrt3/\tau$, os três estágios dão defasagem de $180^\circ$ e atenuação de $8$ vezes. & dedução no capítulo & teoria \\
7 & O erro de nível é $1/(1+K)$, por isso esse regulador não tem precisão melhor que uns $10\,\%$ (proposição~\ref{prop:cascade}) & Consequência da linha~6 para realimentação proporcional. O eixo real tem realimentação em vários estágios e em várias faixas de tempo (linha~5), de modo que o limite se refere ao modelo~\eqref{eq:cascade} e não foi medido. & dedução no capítulo & teoria \\
8 & Com $K=4$ e $7$ o nível se estabiliza; com $K=12$, pulsações regulares com período de $3.7$--$3.9\,\tau$ (fig.~\ref{fig:cascade}) & Cálculo com \texttt{models/\allowbreak chem\_\allowbreak models.py}: com $K=12$, períodos sucessivos de $3.9$, $3.7$, $3.8$, $3.7\,\tau$; com $K=4$, amplitude de $0.003$ no fim do cálculo. & \texttt{models/\allowbreak chem\_\allowbreak models.py} & modelo \\
9 & O hormônio do estresse sai em pulsos cerca de uma vez por hora; os pulsos nascem da própria realimentação, sem gerador separado & Ratos, infusão constante de hormônio liberador: ACTH e corticosterona oscilam com frequência fisiológica, quase horária; uma entrada rítmica do hipotálamo não é necessária para os pulsos. Um modelo hipófise--adrenal com realimentação atrasada produz essas oscilações. & \cite{Walker2012glc, Walker2010} & hipótese \\
10 & O destinatário lê a frequência das porções, não o nível & Macacos sem liberação própria de hormônio liberador de gonadotrofinas: a infusão constante não restabelece a liberação dos hormônios da hipófise, as porções de hora em hora restabelecem; a volta à infusão constante suprime de novo a resposta. A fadiga do receptor como filtro passa-altas é uma interpretação. & \cite{Belchetz1978} & medido \\
11 & Frequências diferentes de porções agem de modo diferente sobre células diferentes da mesma glândula & Os mesmos macacos: aumentar as porções para $2$--$5$ por hora reduz os dois hormônios da hipófise; reduzi-las a uma a cada $3$~horas diminui um (LH) e aumenta o outro (FSH), de modo que a razão entre eles é definida pela frequência. & \cite{Wildt1981gnrh} & medido \\
12 & O ritmo diário nasce de uma realimentação negativa intracelular com atraso de algumas horas & Revisão: as proteínas dos genes do relógio suprimem, com atraso, a transcrição dos próprios genes; o laço de transcrição e tradução dá um período de cerca de um dia. & \cite{TakahashiJS2017} & medido \\
13 & O gerador principal é um grupo de dezenas de milhares de neurônios que sincroniza o resto do corpo & Revisão: o núcleo supraquiasmático é o gerador diário principal; seus neurônios isolados em cultura são geradores autônomos, e a rede os sincroniza; no camundongo, cerca de $10^4$ neurônios de cada lado. & \cite{Welsh2010scn} & medido \\
14 & O período próprio no ser humano é de $24.18$~horas & Pessoas sob luz fraca, com horário desvinculado do dia: o período dos ritmos de melatonina, temperatura e cortisol foi em média de $24.18$~horas em jovens e idosos, com dispersão estreita. & \cite{Czeisler1999} & medido \\
15 & A luz acerta o gerador como uma malha de captura de fase~\eqref{eq:pll}; é preciso um deslocamento de cerca de $11$~min por dia & Pessoas, um único pulso de luz intensa de $6.7$~horas em diferentes horas do dia: curva de resposta de fase do tipo~1 com amplitude de $5$~horas: atraso antes do mínimo da temperatura corporal, avanço depois. A equação~\eqref{eq:pll} é o modelo padrão; $11$~min $=0.18$~hora é aritmética. & \cite{Khalsa2003prc} & medido \\
16 & Sem luz não há ajuste, e o ritmo segue o período próprio & Pessoas totalmente cegas, sem percepção de luz, vivendo em sociedade normal: em cerca de metade delas os ritmos de melatonina e cortisol seguem um período próprio, que não coincide com o dia. & \cite{Sack1992blind} & medido \\
\bottomrule
\end{longtable}}
